\section{Một cuộc kiểm đếm}
\label{sec:ch12-mot-cuoc-kiem-dem}

\index{kiểm đếm tài liệu XAI|(}%
Mục trước đặt một câu hỏi thực nghiệm: khối thí nghiệm có người mà cách đánh
giá theo chức năng vay mượn có lớn tới đâu. Suh, Hurley, Smith và Siu, bốn tác
giả ở MIT Lincoln Laboratory, trả lời câu ấy bằng cách đi đếm, và nhan đề bài
báo của họ đã là câu trả lời: \enquote{Fewer Than 1\% of Explainable AI Papers
Validate Explainability with Humans}~\autocite{p24humangap}.

Trước khi đọc con số thì phải biết bài báo này là loại công trình gì, vì điều
đó quyết định con số nói được gì. Đây là một công trình bảy trang trong nhánh
Extended Abstracts của hội nghị CHI 2025, và các tác giả gọi nó là một báo cáo
sơ bộ, \enquote{late-breaking work}, ngay ở câu đầu tiên của phần tóm tắt; mục
cuối bài mang tiêu đề \enquote{Beyond this Late-Breaking Work} và mô tả công
trình như một khúc dạo cho nghiên cứu về sau~\autocite{p24humangap}. Nó không
phải một bài báo đầy đủ và không tự nhận là thế.

\index{kiểm đếm tài liệu XAI!đối tượng được đo}%
Điều nó đo cũng không phải là mô hình. Không có mô hình nào được huấn luyện,
không có bộ dữ liệu dự đoán nào, không chỉ số nào được tính. Vật được đo là
chính khối tài liệu XAI: bao nhiêu bài tuyên bố rằng phương pháp của mình giải
thích được cho người, và bao nhiêu bài trong số đó đem lời tuyên bố ấy ra thử
với người thật. Một cuộc kiểm đếm như vậy đứng hay đổ ở thủ tục tìm và thủ tục
chấm, nên hai thủ tục ấy đáng đọc kỹ.

\index{kiểm đếm tài liệu XAI!thủ tục tìm}%
Thủ tục tìm chạy trên cơ sở dữ liệu Scopus, gồm cả tài liệu đã xuất bản lẫn
bản tiền ấn phẩm. Các tác giả làm việc với một thủ thư chuyên nghiệp trong ba
tháng để dựng và tinh chỉnh bộ từ khóa, lặp lại nhiều vòng, và họ in nguyên
văn hai chuỗi truy vấn Scopus vào bảng 1 của bài để người khác dựng
lại được~\autocite{p24humangap}. Truy vấn thứ nhất gom mọi bài mang từ khóa về
\tn{khả năng giải thích}{explainability} và
\tn{khả năng diễn giải}{interpretability}; truy vấn thứ hai lấy đúng tập ấy
rồi thêm ràng buộc về thử nghiệm với người dùng cuối, bằng những cụm như
\enquote{human-subject}, \enquote{end-user}, \enquote{experiment},
\enquote{feedback} và \enquote{interview}.

Ba loại tài liệu bị loại ngay trong truy vấn, và cả ba đều để lại dấu trên kết
quả. Bài rà tổng quan bị loại bằng điều kiện tiêu đề chứa \enquote{review},
\enquote{survey} hoặc \enquote{overview} và bằng loại tài liệu; bài không
tiếng Anh bị loại bằng điều kiện ngôn ngữ; bài hội thảo nhỏ bị loại theo mô tả
của chính các tác giả ở mục bàn về hướng tiếp theo. Đó là ba giới hạn mà bài
báo tự nêu, chứ không phải ba điều sách phát hiện ra.

\index{kiểm đếm tài liệu XAI!thủ tục chấm}%
Thủ tục chấm thì mượn từ chỗ khác, và chỗ mượn đáng ghi lại. Tiêu chí chấm lấy
từ Miller và cộng sự: một bài được 1 điểm nếu chính nó kiểm chứng lời tuyên bố
về khả năng giải thích với người, và 0 điểm nếu không. Mỗi bài trong tập cuối
được ba người chấm, và chỗ bất đồng được giải quyết bằng thảo
luận~\autocite{p24humangap}.

Có một chỗ trong thủ tục chấm mà bài báo thu hẹp phạm vi có chủ ý, và chỗ ấy
đổi nghĩa của con số cuối cùng. Miller và cộng sự có hai tiêu chí, không phải
một: bên cạnh tiêu chí kiểm chứng vừa nêu còn một tiêu chí thứ hai, chấm theo
việc bài báo có dẫn hoặc có triển khai công trình khoa học xã hội hay không.
Suh và cộng sự không dùng tiêu chí thứ hai, và họ nói lý do: nó được đặt ra
năm 2017 nên \enquote{may be outdated}, và họ giữ phạm vi ở \enquote{only the
use of humans in a behavioral study to evaluate the XAI
method}~\autocite{p24humangap}. Vậy nên cái được đếm là sự có mặt của một thí
nghiệm hành vi với người, không phải chất lượng của thí nghiệm ấy, và cũng
không phải mức độ bài báo dựa vào hiểu biết đã có về con người.

Bài báo còn liệt kê bốn tình huống mà một bài bị chấm 0 dù thoạt nhìn có vẻ
đáng được điểm, và bốn tình huống ấy vẽ ra ranh giới của phép đếm rõ hơn mọi
định nghĩa: bài tự nhận là đã diễn giải được nhờ chính thiết kế của mô hình
nên cho rằng mình không cần thí nghiệm với người; bài đóng góp một khung hoặc
một cách đánh giá tự nhận là giải thích được nhưng không đem khung ấy thử với
người; bài dùng dữ
liệu do người sinh ra trong lúc huấn luyện nhưng không đánh giá với người; và
bài so kết quả của mình với nhãn do người gán mà không có bước kiểm chứng nào
với người~\autocite{p24humangap}. Tình huống thứ tư là tình huống dễ nhầm nhất
trong bốn, vì nhãn do người gán trông rất giống một phép kiểm chứng có người,
và mục~\ref{sec:ch12-nghien-cuu-khong-co-nguoi} quay lại đúng chỗ ấy.
